\documentclass[12pt,oneside]{book}

%============================================================
% PACKAGES
%============================================================

%--------------------
% Page + font
%--------------------
\usepackage[margin=1in]{geometry}
\usepackage[T1]{fontenc}
\usepackage[utf8]{inputenc}
\usepackage{lmodern}
\usepackage{microtype}

%--------------------
% Math
%--------------------
\usepackage{amsmath, amssymb, amsfonts, amsthm, mathtools}
\usepackage{bm}
\usepackage{bbm}
\usepackage{mathrsfs}
\usepackage{dsfont}

%--------------------
% Tables + figures
%--------------------
\usepackage{graphicx}
\usepackage{float}
\usepackage{booktabs}
\usepackage{array}
\usepackage{xltabular}
\usepackage{multirow}
\usepackage{caption}
\usepackage{subcaption}

%--------------------
% Lists
%--------------------
\usepackage{enumitem}
\setlist[itemize]{topsep=2pt,itemsep=2pt}
\setlist[enumerate]{topsep=2pt,itemsep=2pt}

%--------------------
% Colors + hyperlinks
%--------------------
\usepackage[dvipsnames]{xcolor}
\usepackage[
    colorlinks=true,
    linkcolor=MidnightBlue,
    citecolor=MidnightBlue,
    urlcolor=MidnightBlue
]{hyperref}
\usepackage[nameinlink,capitalise]{cleveref}

%--------------------
% Algorithms
%--------------------
\usepackage{algorithm}
\usepackage{algpseudocode}

%--------------------
% Code blocks
%--------------------
\usepackage{listings}

\lstdefinestyle{codestyle}{
    basicstyle=\ttfamily\small,
    breaklines=true,
    frame=single,
    numbers=left,
    numberstyle=\tiny,
    keywordstyle=\color{MidnightBlue},
    commentstyle=\color{ForestGreen},
    stringstyle=\color{BrickRed},
    showstringspaces=false,
    tabsize=4
}

\lstset{style=codestyle}

%--------------------
% Section styling
%--------------------
\usepackage{titlesec}
\usepackage{tocloft}

\setcounter{tocdepth}{2}
\setcounter{secnumdepth}{3}

%============================================================
% THEOREMS
%============================================================

\numberwithin{equation}{section}

\theoremstyle{definition}
\newtheorem{definition}{Definition}[section]
\newtheorem{example}{Example}[section]
\newtheorem{exercise}{Exercise}[section]
\newtheorem{remark}{Remark}[section]

\theoremstyle{plain}
\newtheorem{theorem}{Theorem}[section]
\newtheorem{lemma}{Lemma}[section]
\newtheorem{proposition}{Proposition}[section]
\newtheorem{corollary}{Corollary}[section]

\theoremstyle{remark}
\newtheorem*{note}{Note}

%============================================================
% CUSTOM COMMANDS
%============================================================

% Sets
\newcommand{\R}{\mathbb{R}}
\newcommand{\N}{\mathbb{N}}
\newcommand{\Z}{\mathbb{Z}}
\newcommand{\Q}{\mathbb{Q}}
\newcommand{\C}{\mathbb{C}}

% Probability
\newcommand{\E}{\mathbb{E}}
\newcommand{\Var}{\mathrm{Var}}
\newcommand{\Cov}{\mathrm{Cov}}
\newcommand{\Corr}{\mathrm{Corr}}
\newcommand{\Prob}{\mathbb{P}}

\newcommand{\ind}{\mathbbm{1}}
\newcommand{\iid}{\overset{\text{iid}}{\sim}}
\newcommand{\convd}{\overset{d}{\to}}
\newcommand{\convp}{\overset{p}{\to}}
\newcommand{\as}{\overset{a.s.}{\to}}

% Operators
\newcommand{\given}{\,\middle|\,}
\newcommand{\argmin}{\operatorname*{arg\,min}}
\newcommand{\argmax}{\operatorname*{arg\,max}}

% Linear algebra / analysis
\newcommand{\norm}[1]{\left\lVert #1 \right\rVert}
\newcommand{\abs}[1]{\left\lvert #1 \right\rvert}
\newcommand{\inner}[2]{\left\langle #1, #2 \right\rangle}

\newcommand{\dd}{\,\mathrm{d}}
\newcommand{\grad}{\nabla}
\newcommand{\Hess}{\nabla^2}

\newcommand{\tr}{\operatorname{tr}}
\newcommand{\rank}{\operatorname{rank}}
\newcommand{\diag}{\operatorname{diag}}

% Distributions
\newcommand{\Normal}{\mathcal{N}}
\newcommand{\Bern}{\mathrm{Bernoulli}}
\newcommand{\Bin}{\mathrm{Binomial}}
\newcommand{\Pois}{\mathrm{Poisson}}
\newcommand{\Gam}{\mathrm{Gamma}}
\newcommand{\BetaDist}{\mathrm{Beta}}
\newcommand{\Exp}{\mathrm{Exponential}}
\newcommand{\Dir}{\mathrm{Dirichlet}}
\newcommand{\InvGam}{\mathrm{Inv\text{-}Gamma}}

% Title block
\title{EN.553.632 Bayesian Statistics}
\author{Jonathan Y. Ma}
\date{June 2026}

\begin{document}

\maketitle
\tableofcontents
\chapter{Introduction and Probability Review}

\section{Bayesian vs Frequentist Perspectives}

\begin{definition}
Probability can be interpreted in two fundamentally different ways:
\begin{itemize}
    \item \textbf{Frequentist:} Probability is the long-run relative frequency of an event.
    \item \textbf{Bayesian:} Probability is a degree of belief about an uncertain quantity.
\end{itemize}
\end{definition}

In the Bayesian framework, unknown parameters are treated as random variables. This contrasts with the frequentist view, where parameters are fixed but unknown.

\begin{example}
Consider a coin flipped $n=10$ times with $7$ heads observed.

\begin{itemize}
    \item Frequentist estimate: $\hat{\theta} = 0.7$
    \item Bayesian approach: place a prior $\theta \sim \BetaDist(\alpha,\beta)$
\end{itemize}

After observing $y=7$ heads, the posterior becomes:
$$
\theta \mid y \sim \BetaDist(\alpha + y, \beta + n - y)
$$
\end{example}

\section{Bayes' Theorem}

\begin{theorem}[Bayes' Rule]
For parameter $\theta$ and data $y$,
$$
p(\theta \mid y) = \frac{p(y \mid \theta)p(\theta)}{p(y)}
$$
where
$$
p(y) = \int p(y \mid \theta)p(\theta)\dd \theta
$$
\end{theorem}

\begin{remark}
The posterior is often written up to proportionality:
$$
p(\theta \mid y) \propto p(y \mid \theta)p(\theta)
$$
\end{remark}

\section{Bernoulli and Binomial Models}

\subsection{Bernoulli Model}

\begin{definition}
A Bernoulli random variable satisfies:
$$
Y \sim \Bern(\theta), \quad P(Y=1)=\theta
$$
\end{definition}

\subsection{Binomial Model}

\begin{definition}
For $n$ independent Bernoulli trials,
$$
Y \sim \Bin(n,\theta), \quad P(Y=y|\theta) = \binom{n}{y}\theta^y(1-\theta)^{n-y}
$$
\end{definition}

\subsection{Likelihood}

Given observations $y_1,\dots,y_n$, the likelihood is:
$$
L(\theta) = \prod_{i=1}^n \theta^{y_i}(1-\theta)^{1-y_i}
= \theta^{\sum y_i}(1-\theta)^{n-\sum y_i}
$$

\subsection{Maximum Likelihood Estimation}

\begin{proposition}
The MLE for $\theta$ in the Bernoulli/Binomial model is:
$$
\hat{\theta}_{\text{MLE}} = \bar{y}
$$
\end{proposition}

\begin{proof}
The log-likelihood is:
$$
\ell(\theta) = \sum y_i \log \theta + (n-\sum y_i)\log(1-\theta)
$$

Differentiate:
$$
\frac{d\ell}{d\theta} = \frac{\sum y_i}{\theta} - \frac{n-\sum y_i}{1-\theta}
$$

Set equal to zero:
$$
\frac{\sum y_i}{\theta} = \frac{n-\sum y_i}{1-\theta}
\Rightarrow \theta = \frac{1}{n}\sum y_i
$$
\end{proof}

\section{Conjugate Priors: Beta-Binomial Model}

\begin{definition}
A prior is conjugate if the posterior belongs to the same family as the prior.
\end{definition}

\begin{proposition}
If
$$
\theta \sim \BetaDist(\alpha,\beta), \quad Y \mid \theta \sim \Bin(n,\theta)
$$
then
$$
\theta \mid y \sim \BetaDist(\alpha + y, \beta + n - y)
$$
\end{proposition}

\begin{proof}
Posterior:
$$
p(\theta \mid y) \propto \theta^y(1-\theta)^{n-y} \cdot \theta^{\alpha-1}(1-\theta)^{\beta-1}
$$

Combine exponents:
$$
p(\theta \mid y) \propto \theta^{\alpha+y-1}(1-\theta)^{\beta+n-y-1}
$$

which is the kernel of a Beta distribution.
\end{proof}

\subsection{Posterior Mean}

$$
\E[\theta \mid y] = \frac{\alpha + y}{\alpha + \beta + n}
$$

\begin{remark}
This can be interpreted as a weighted average:
$$
\E[\theta \mid y] = w \cdot \frac{\alpha}{\alpha+\beta} + (1-w)\cdot \frac{y}{n}
$$
\end{remark}

\section{Credible Intervals and HPD Regions}

\begin{definition}
An interval $(a,b)$ is a $(1-\alpha)$ credible interval if
$$
\Prob(a < \theta < b \mid y) = 1 - \alpha
$$
\end{definition}

\begin{definition}
A highest posterior density (HPD) region $A$ satisfies:
\begin{itemize}
    \item $\Prob(\theta \in A \mid y) = 1-\alpha$
    \item $p(\theta_1 \mid y) > p(\theta_2 \mid y)$ for $\theta_1 \in A$, $\theta_2 \notin A$
\end{itemize}
\end{definition}

\section{Posterior Predictive Distribution}

\begin{proposition}
For a new observation $Z$,
$$
p(Z=1 \mid y) = \E[\theta \mid y]
$$
\end{proposition}

\begin{proof}
$$
p(Z=1 \mid y) = \int p(Z=1 \mid \theta)p(\theta \mid y)\dd \theta
= \int \theta p(\theta \mid y)\dd \theta = \E[\theta \mid y]
$$
\end{proof}

\section{Exponential Family}

\begin{definition}
A distribution belongs to the exponential family if it can be written as:
$$
p(y \mid \theta) = f(y)g(\theta)\exp\{\phi(\theta)u(y)\}
$$
\end{definition}

\begin{remark}
Many common distributions (Normal, Binomial, Poisson) belong to this family, which enables conjugate prior constructions.
\end{remark}

\section{Prior Specification}

\subsection{Subjective Priors}

Constructed from expert knowledge.

\subsection{Non-informative Priors}

Designed to have minimal influence on the posterior.

\begin{example}
A flat prior:
$$
\pi(\theta) \propto 1
$$
\end{example}

\subsection{Jeffreys Prior}

\begin{definition}
$$
\pi(\theta) \propto \sqrt{I(\theta)}, \quad I(\theta) = -\E\left[\frac{\partial^2}{\partial \theta^2}\log p(y \mid \theta)\right]
$$
\end{definition}

\begin{remark}
Jeffreys priors are invariant under reparameterization.
\end{remark}

\section{Poisson-Gamma Model}

\subsection{Model}

$$
Y_i \iid \Pois(\theta)
$$

Likelihood:
$$
L(\theta) \propto \theta^{\sum y_i} e^{-n\theta}
$$

\subsection{Prior}

$$
\theta \sim \Gam(a,b)
$$

\subsection{Posterior}

\begin{proposition}
$$
\theta \mid y \sim \Gam(a + \sum y_i, b + n)
$$
\end{proposition}

\section{Preview: Computational Methods}

In many models, posterior distributions cannot be computed analytically. This motivates:

\begin{itemize}
    \item Monte Carlo integration
    \item Importance sampling
    \item Markov Chain Monte Carlo (MCMC)
\end{itemize}

These methods will be developed in later chapters.

\section{Common Distributions and Identities}

\subsection{Beta and Gamma Functions}

\begin{definition}
The Gamma function is defined as:
$$
\Gamma(a) = \int_0^\infty x^{a-1} e^{-x} \dd x
$$
and satisfies:
$$
\Gamma(n) = (n-1)! \quad \text{for integer } n
$$
\end{definition}

\begin{definition}
The Beta distribution has density:
$$
f(\theta) = \frac{\Gamma(a+b)}{\Gamma(a)\Gamma(b)} \theta^{a-1}(1-\theta)^{b-1}
$$
\end{definition}

\begin{proposition}
For $\theta \sim \BetaDist(a,b)$:
$$
\E[\theta] = \frac{a}{a+b}, \quad 
\Var(\theta) = \frac{ab}{(a+b)^2(a+b+1)}
$$
\end{proposition}

\subsection{Gamma Distribution}

\begin{definition}
A Gamma distribution with shape $a$ and rate $b$:
$$
f(\theta) = \frac{b^a}{\Gamma(a)} \theta^{a-1} e^{-b\theta}
$$
\end{definition}

\begin{proposition}
$$
\E[\theta] = \frac{a}{b}, \quad \Var(\theta) = \frac{a}{b^2}
$$
\end{proposition}

\begin{remark}
The Gamma distribution is conjugate to the Poisson likelihood.
\end{remark}

\section{Sufficiency and Likelihood Structure}

\begin{definition}
A statistic $T(y)$ is sufficient for $\theta$ if the likelihood can be written as:
$$
p(y \mid \theta) = g(T(y), \theta) h(y)
$$
\end{definition}

\begin{proposition}
For:
\begin{itemize}
    \item Binomial model: sufficient statistic is $y = \sum y_i$
    \item Poisson model: sufficient statistic is $\sum y_i$
\end{itemize}
\end{proposition}

\begin{remark}
This reduction is critical in Bayesian updating since posterior distributions depend only on sufficient statistics.
\end{remark}

\section{Posterior as a Weighted Average}

A key Bayesian insight is that posterior estimates combine prior and data:

\begin{proposition}
For the Beta-Binomial model:
$$
\E[\theta \mid y] = \frac{a+y}{a+b+n}
$$
can be written as:
$$
\E[\theta \mid y] =
\frac{a+b}{a+b+n} \cdot \frac{a}{a+b}
+
\frac{n}{a+b+n} \cdot \frac{y}{n}
$$
\end{proposition}

\begin{remark}
This shows:
\begin{itemize}
    \item Prior contributes as pseudo-observations
    \item Data influence increases with $n$
\end{itemize}
\end{remark}

\section{Bias, Variance, and MSE}

\begin{definition}
Bias of an estimator $\hat{\theta}$:
$$
\text{Bias} = \E[\hat{\theta}] - \theta
$$
\end{definition}

\begin{definition}
Mean Squared Error:
$$
\text{MSE}(\hat{\theta}) = \Var(\hat{\theta}) + \text{Bias}^2
$$
\end{definition}

\begin{remark}
Bayesian estimators are often biased but achieve lower MSE through shrinkage.
\end{remark}

\section{Poisson-Gamma Model (Detailed)}

\subsection{Likelihood}

$$
L(\theta) = \prod_{i=1}^n \frac{\theta^{y_i} e^{-\theta}}{y_i!}
\propto \theta^{\sum y_i} e^{-n\theta}
$$

\subsection{Posterior Derivation}

\begin{proof}
With prior $\theta \sim \Gam(a,b)$:
$$
p(\theta \mid y) \propto \theta^{\sum y_i} e^{-n\theta} \cdot \theta^{a-1} e^{-b\theta}
$$

$$
= \theta^{a+\sum y_i -1} e^{-(b+n)\theta}
$$

which is $\Gam(a+\sum y_i, b+n)$.
\end{proof}

\subsection{Posterior Mean Interpretation}

$$
\E[\theta \mid y] =
\frac{b}{b+n} \cdot \frac{a}{b}
+
\frac{n}{b+n} \cdot \frac{\sum y_i}{n}
$$

\begin{remark}
Again, this is a weighted average of prior mean and sample mean.
\end{remark}

\section{Posterior Predictive Distribution}

\begin{definition}
The posterior predictive distribution is:
$$
p(\tilde{y} \mid y) = \int p(\tilde{y} \mid \theta)p(\theta \mid y)\dd \theta
$$
\end{definition}

\begin{remark}
This integrates out parameter uncertainty, a key advantage over plug-in estimators.
\end{remark}

\section{Exponential Family and Conjugacy}

\begin{definition}
A one-parameter exponential family has the form:
$$
p(y \mid \phi) = h(y)c(\phi)\exp\{\phi t(y)\}
$$
\end{definition}

\begin{proposition}
For exponential families:
\begin{itemize}
    \item Sufficient statistics appear naturally as $t(y)$
    \item Conjugate priors exist due to exponential structure
\end{itemize}
\end{proposition}

\begin{remark}
Examples:
\begin{itemize}
    \item Binomial $\leftrightarrow$ Beta
    \item Poisson $\leftrightarrow$ Gamma
    \item Normal $\leftrightarrow$ Normal / Inverse-Gamma
\end{itemize}
\end{remark}

\section{Interpretation of Bayesian Updating}

\begin{remark}
Bayesian inference provides:
\begin{itemize}
    \item Full uncertainty quantification (posterior distribution)
    \item Direct probabilistic statements about parameters
    \item Natural predictive distributions
\end{itemize}
\end{remark}

\begin{example}
In rare-event estimation, even if $y=0$, the posterior remains well-defined:
$$
\theta \mid y \sim \BetaDist(a, b+n)
$$
which avoids degenerate estimates like $\hat{\theta}=0$.
\end{example}

\chapter{The Binomial Model}

\section{Model Setup}

\begin{definition}
Let $Y \mid \theta \sim \Bin(n,\theta)$ with probability mass function:
$$
p(y \mid \theta) = \binom{n}{y}\theta^y(1-\theta)^{n-y}, \quad \theta \in [0,1]
$$
\end{definition}

\begin{remark}
If $Y_i \iid \Bern(\theta)$, then $Y = \sum_{i=1}^n Y_i \sim \Bin(n,\theta)$. The sufficient statistic is $y = \sum Y_i$.
\end{remark}

\section{Likelihood and Maximum Likelihood Estimation}

Ignoring constants:
$$
L(\theta) \propto \theta^y(1-\theta)^{n-y}
$$

\begin{proposition}
The log-likelihood is:
$$
\ell(\theta) = y\log \theta + (n-y)\log(1-\theta)
$$
\end{proposition}

\begin{proposition}
The MLE for $\theta$ is:
$$
\hat{\theta} = \frac{y}{n}
$$
\end{proposition}

\begin{proof}
Differentiate:
$$
\frac{d\ell}{d\theta} = \frac{y}{\theta} - \frac{n-y}{1-\theta}
$$

Set equal to zero:
$$
\frac{y}{\theta} = \frac{n-y}{1-\theta}
\Rightarrow \theta = \frac{y}{n}
$$
\end{proof}

\section{Beta Prior}

\begin{definition}
Assume prior:
$$
\theta \sim \BetaDist(a,b), \quad \pi(\theta) \propto \theta^{a-1}(1-\theta)^{b-1}
$$
\end{definition}

\begin{remark}
The parameters $(a,b)$ can be interpreted as prior pseudo-counts:
$$
a-1 \text{ successes}, \quad b-1 \text{ failures}
$$
\end{remark}

\section{Posterior Distribution}

\begin{theorem}
If $\theta \sim \BetaDist(a,b)$ and $Y \mid \theta \sim \Bin(n,\theta)$, then:
$$
\theta \mid y \sim \BetaDist(a+y, b+n-y)
$$
\end{theorem}

\begin{proof}
$$
p(\theta \mid y) \propto \theta^y(1-\theta)^{n-y} \cdot \theta^{a-1}(1-\theta)^{b-1}
$$

$$
= \theta^{a+y-1}(1-\theta)^{b+n-y-1}
$$
\end{proof}

\section{Posterior Summaries}

\subsection{Mean}

$$
\E[\theta \mid y] = \frac{a+y}{a+b+n}
$$

\begin{remark}
This can be written as:
$$
\E[\theta \mid y]
=
\frac{a+b}{a+b+n}\cdot \frac{a}{a+b}
+
\frac{n}{a+b+n}\cdot \frac{y}{n}
$$
\end{remark}

\subsection{Mode}

$$
\text{Mode}(\theta \mid y) =
\frac{a+y-1}{a+b+n-2}
$$

\subsection{Variance}

$$
\Var(\theta \mid y)
=
\frac{(a+y)(b+n-y)}{(a+b+n)^2(a+b+n+1)}
$$

\section{Posterior Predictive Distribution}

\begin{definition}
For a new observation $\tilde{Y}$:
$$
p(\tilde{Y}=1 \mid y)
=
\int \theta p(\theta \mid y)\dd \theta
$$
\end{definition}

\begin{proposition}
$$
p(\tilde{Y}=1 \mid y) = \E[\theta \mid y] = \frac{a+y}{a+b+n}
$$
\end{proposition}

\section{Credible Intervals}

\begin{definition}
An interval $(l,u)$ is a $(1-\alpha)$ credible interval if:
$$
\Prob(l < \theta < u \mid y) = 1-\alpha
$$
\end{definition}

\section{Highest Posterior Density Regions}

\begin{definition}
A set $A$ is a HPD region if:
\begin{itemize}
\item $\Prob(\theta \in A \mid y) = 1-\alpha$
\item $p(\theta_1 \mid y) > p(\theta_2 \mid y)$ for $\theta_1 \in A$, $\theta_2 \notin A$
\end{itemize}
\end{definition}

\section{Shrinkage Interpretation}

\begin{proposition}
The posterior mean can be written as:
$$
\E[\theta \mid y]
=
w \cdot \frac{a}{a+b}
+
(1-w)\cdot \frac{y}{n}
$$

where:
$$
w = \frac{a+b}{a+b+n}
$$
\end{proposition}

\begin{remark}
\begin{itemize}
\item Small $n$: prior dominates
\item Large $n$: data dominates
\end{itemize}
\end{remark}

\section{Rare Event Example}

\begin{example}
Suppose:
$$
n=20, \quad y=0, \quad \theta \sim \BetaDist(2,20)
$$

Then:
$$
\theta \mid y \sim \BetaDist(2,40)
$$
\end{example}

\begin{remark}
The MLE gives $\hat{\theta}=0$, while the Bayesian estimator remains strictly positive.
\end{remark}

\section{Frequentist vs Bayesian Interpretation}

\begin{remark}
Frequentist interval:
$$
\Prob(l(Y) < \theta < u(Y) \mid \theta) = 1-\alpha
$$

Bayesian interval:
$$
\Prob(l < \theta < u \mid y) = 1-\alpha
$$
\end{remark}

\chapter{Poisson Model and Exponential Families}

\section{Poisson Model}

\begin{definition}
A random variable $Y$ follows a Poisson distribution with rate $\theta > 0$ if:
$$
Y \mid \theta \sim \Pois(\theta), \quad
p(y \mid \theta) = \frac{\theta^y e^{-\theta}}{y!}, \quad y \in \N
$$
\end{definition}

\begin{proposition}
For $Y \sim \Pois(\theta)$:
$$
\E[Y] = \theta, \quad \Var(Y) = \theta
$$
\end{proposition}

\begin{remark}
The equality of mean and variance is a defining feature of the Poisson model. Large deviations between sample mean and variance indicate model misspecification.
\end{remark}

\section{Likelihood and Sufficiency}

Let $Y_1,\dots,Y_n \iid \Pois(\theta)$. The joint likelihood is:
$$
p(y_1,\dots,y_n \mid \theta)
= \prod_{i=1}^n \frac{\theta^{y_i} e^{-\theta}}{y_i!}
= c(y)\,\theta^{\sum y_i} e^{-n\theta}
$$

\begin{proposition}
The sufficient statistic is:
$$
T(y) = \sum_{i=1}^n y_i
$$
\end{proposition}

\begin{remark}
All information about $\theta$ is contained in the total count $\sum y_i$.
\end{remark}

\section{Maximum Likelihood Estimation}

\begin{proposition}
The MLE for $\theta$ is:
$$
\hat{\theta} = \bar{y} = \frac{1}{n}\sum_{i=1}^n y_i
$$
\end{proposition}

\begin{proof}
Log-likelihood:
$$
\ell(\theta) = \sum y_i \log \theta - n\theta
$$

Differentiate:
$$
\frac{d\ell}{d\theta} = \frac{\sum y_i}{\theta} - n
$$

Set to zero:
$$
\theta = \frac{1}{n}\sum y_i
$$
\end{proof}

\section{Gamma Prior}

\begin{definition}
Assume prior:
$$
\theta \sim \Gam(a,b), \quad
\pi(\theta) = \frac{b^a}{\Gamma(a)} \theta^{a-1} e^{-b\theta}
$$
\end{definition}

\begin{proposition}
$$
\E[\theta] = \frac{a}{b}, \quad \Var(\theta) = \frac{a}{b^2}
$$
\end{proposition}

\section{Posterior Distribution}

\begin{theorem}
If $\theta \sim \Gam(a,b)$ and $Y_i \iid \Pois(\theta)$, then:
$$
\theta \mid y \sim \Gam\left(a + \sum y_i,\; b + n\right)
$$
\end{theorem}

\begin{proof}
$$
p(\theta \mid y)
\propto \theta^{\sum y_i} e^{-n\theta} \cdot \theta^{a-1} e^{-b\theta}
$$

$$
= \theta^{a+\sum y_i -1} e^{-(b+n)\theta}
$$
\end{proof}

\section{Posterior Summaries}

\subsection{Mean}

$$
\E[\theta \mid y]
= \frac{a + \sum y_i}{b + n}
$$

\begin{remark}
This can be written as:
$$
\E[\theta \mid y]
=
\frac{b}{b+n}\cdot \frac{a}{b}
+
\frac{n}{b+n}\cdot \frac{\sum y_i}{n}
$$
\end{remark}

\subsection{Variance}

$$
\Var(\theta \mid y)
=
\frac{a + \sum y_i}{(b+n)^2}
$$

\section{Posterior Predictive Distribution}

\begin{definition}
For a new observation $\tilde{Y}$:
$$
p(\tilde{y} \mid y)
=
\int p(\tilde{y} \mid \theta)p(\theta \mid y)\dd \theta
$$
\end{definition}

\begin{remark}
This results in a negative binomial distribution, reflecting overdispersion relative to the Poisson model.
\end{remark}

\section{Example: Birth Rate Comparison}

\begin{example}
Suppose:
\begin{itemize}
\item Group 1: $n_1=111$, $\sum y_i = 217$
\item Group 2: $n_2=44$, $\sum y_i = 66$
\item Prior: $\theta \sim \Gam(2,1)$
\end{itemize}

Posteriors:
$$
\theta_1 \mid y \sim \Gam(219,112), \quad
\theta_2 \mid y \sim \Gam(68,45)
$$
\end{example}

\section{Exponential Family Representation}

\begin{definition}
A distribution belongs to the exponential family if:
$$
p(y \mid \phi) = h(y)c(\phi)\exp\{\phi t(y)\}
$$
\end{definition}

\subsection{Poisson as an Exponential Family}

\begin{proposition}
The Poisson distribution can be written as:
$$
p(y \mid \theta)
= \frac{1}{y!}\exp\{y\log \theta - \theta\}
$$
\end{proposition}

\begin{remark}
This corresponds to:
$$
t(y) = y, \quad \phi = \log \theta
$$
\end{remark}

\section{General Exponential Family Structure}

\begin{definition}
A one-parameter exponential family has the form:
$$
p(y \mid \phi) = h(y)c(\phi)\exp\{\phi t(y)\}
$$
\end{definition}

\begin{proposition}
For exponential families:
\begin{itemize}
\item $t(y)$ is a sufficient statistic
\item Conjugate priors exist due to exponential structure
\end{itemize}
\end{proposition}

\begin{remark}
Common examples:
\begin{itemize}
\item Binomial $\leftrightarrow$ Beta
\item Poisson $\leftrightarrow$ Gamma
\item Normal $\leftrightarrow$ Normal / Inverse-Gamma
\end{itemize}
\end{remark}

\section{Interpretation and Model Limitations}

\begin{remark}
The Poisson model assumes:
\begin{itemize}
\item independence of events
\item constant rate $\theta$
\item equal mean and variance
\end{itemize}
\end{remark}

\begin{remark}
If variance exceeds the mean, a negative binomial model may be more appropriate.
\end{remark}



\end{document}